
\textcolor{red}{Não apagar, pois tem formulações referente à punção no perímetro C', que pode ser necessário resgatar.!!!}
%Não apagar!!!!!!!!!
% Para o perímetro $C'$, procedem-se às verificações de punção conforme descrito a seguir. inicialmente, pela equação \eqref{eq:secao_critica_c'} é determinada a distância da seção crítica em relação ao  pilar: 

% \begin{equation}
% sc = \frac{1,00}{2} = 0,50\;m
%     \label{eq:EL16_secao_critica}
% \end{equation}

% Em seguida, com base na equação~\eqref{eq:ke}, define-se o coeficiente de escala para punção:

% \begin{equation}
% k_e = 1 + \sqrt{\frac{20}{0,96 \cdot 100}} = 1,46
% \label{eq:EL16_ke}
% \end{equation}

% \textcolor{red}{pegar essa informação com o professor}
% Para os valores da taxa de armadura, aplica-se a tabela ??? da NBR 6118, que relaciona a taxa de armadura com a resistência característica do concreto. Ao utilizar $f_{ck}$ = 25 MPa, obtem-se os valores de $rho_x = 0,15$ e $rho_y = 0,15$. Assim, de acordo com a equação \eqref{eq:rho}, tem-se:

% \begin{equation}
% rho = \frac{\sqrt{0,15 * 0,15}}{100} = 1,50 \cdot 10^{-3}
%     \label{eq:rho_el16}
% \end{equation}

% A tensão resistente à punção no perímetro $C'$ é definida por meio equação~\eqref{eq:Tensão_resistente_C'}:

% \begin{equation}
%  \tau_{rd,1} = 1000 \cdot \left(0,13 \cdot 1,46\cdot \left( 100 \cdot 1,50\cdot10^{-3} \cdot 25,00 \right) ^{\frac{1}{3}}  + 0,1 \cdot 0,00 \right) = 292,85\;kPa
%  \label{eq:EL16_Tensão_resistente_C'}
% \end{equation}

% O perímetro crítico associado a seção crítica $C'$ é definido pela equação \eqref{eq:perímetro_c'}:

% \begin{equation}
% u_{rd,1} = 2 \cdot \left( 0,25 + 1,20 \right) + 4 \pi \cdot 0,50 = 9,18\;m 
%     \label{eq:EL16_perímetro_c'}
% \end{equation}

% Os valores dos coeficientes $k_x$ e $k_y$, que fornecem a parcela do momento transmitida ao pilar por cisalhamento, são obtidos a partir da tabela 19.2 da NBR 6118, através da relação das dimensões dos pilares. Dessa forma, aplicando as equações \eqref{eq:c1-C2} e \eqref{eq:C2-C1}, tem-se:

% \begin{equation}
% C1\_C2 = \frac{0,25}{1,20} = 0,21
%     \label{eq:EL16_c1-C2}
% \end{equation}

% \begin{equation}
% C2\_C1 = \frac{1,20}{0,25} = 4,80
%     \label{eq:EL16_C2-C1}
% \end{equation}

% Após a analise da referida tabela, definem-se os coeficientes em: $k_x = 0,45$  e $k_y = 0,80$.


% Os módulos de resistência plásticos associados às direções $x$ e $y$ são calculados 
% por meio das equações \eqref{eq:wpx} e \eqref{eq:wpy}:

% \begin{equation}
% w_{px} = \frac{0,25^2}{2} + 0,25\cdot 1,20 + 4 \cdot 1,20 \cdot 0,50 + 16 \cdot 0,5^2 + 2 \cdot \pi \cdot 0,25 \cdot 0,50 = 7,52 \; m^2
% \label{eq:EL16_wpy}
% \end{equation}

% \begin{equation}
% w_{py} = \frac{1,20^2}{2} + 1,20 \cdot 0,25 + 4 \cdot 0,25 \cdot 0,50 + 16 \cdot 0,50^2 + 2 \cdot \pi \cdot 1,2 \cdot 0,50 = 9,29\;m^2
%     \label{eq:EL16_wpx}
% \end{equation}

% Em posse de tais valores, define-se a tensão solicitante de punção no perímetro crítico $C'$ conforme a equação \eqref{eq:Tensão_solicitante_C'}:
 
% \begin{equation}
%  \tau_{sd,1} = \frac{1,4 \cdot 377,30}{9,18 \cdot 0,96} + \frac{1,4\cdot 0,45 \cdot (-1,90)}{7,52 \cdot 0,96} + \frac{1,4 \cdot 0,80 \cdot 4,40}{9,29 \cdot 0,96} = 60,32 \; kPa
%      \label{eq:EL16_Tensão_solicitante_C'}
% \end{equation}

% Por fim, a verificação da restrição associada à punção na seção crítica $C'$ é realizada por meio da Equação~\eqref{eq:restricao_tensao_c'}, resultando em:

% \begin{equation}
% g_{rd,1} = \frac{60,76}{292,86} - 1 = -0,79
%     \label{eq:EL16_grd1}
% \end{equation}

% Nota-se que o valor negativo de $g_{rd,1}$ indica que a condição de resistência à punção no perímetro~$C'$ é plenamente atendida.


% Concluídos os cálculos para a combinação de esforços $c1$, procede-se à repetição da análise para as demais combinações ($c2 $ e $c3$), com o objetivo de verificar as tensões solicitantes no solo e a resistência à punção. Os resultados da verificação de punção na seção crítica $C'$  são apresentados na tabela \ref{tab:puncao_secao_c_linha_p16}, enquanto aqueles referentes à seção crítica $C$ encontram-se na tabela \ref{tab:puncao_secao_c}. Os valores relacionados às tensão solicitantes no solo estão presentes na tabela \ref{tab:tensoes_solo}:

% \begin{table}[H]
% \centering
% \caption{verificação seção $C'$}
% \label{tab:puncao_secao_c_linha_p16}
% \setlength{\tabcolsep}{10pt}
% \begin{tabular}{c c c c c c c c c}
% \toprule

% Elemento & $\tau_{rd,1}$ & $u_{rd,1}$ & $k_x$ & $k_y$ & $w_{p,x}$ & $w_{p,y}$ & $\tau_{sd,1}$ & $g_{rd,1}$ \\
% \midrule
% P16 (c1) & 292,85 & 9,18 & 0,45 & 0,80 & 7,52 & 9,29 & 60,32 & -0,79 \\
% P16 (c2) & 292,85 & 9,18 & 0,45 & 0,80 & 7,52 & 9,29 & 38,15 & -0,87 \\
% P16 (c3) & 292,85 & 9,18 & 0,45 & 0,80 & 7,52 & 9,29 & 62,42 & -0,79 \\
% \bottomrule
% \end{tabular}
% \end{table}


% \begin{table}[H]
% \centering
% \caption{Verificação à punção na seção $C$}
% \label{tab:puncao_secao_c}
% \setlength{\tabcolsep}{10pt}
% \begin{tabular}{lcccccccc}
% \hline
%  Elemento & $\tau_{sd,2}$ & $\tau_{rd,2}$ & $u_{rd,2}$ & $k_e$ & $g_{e,d}$ & $g_{rd,2}$ \\
% \hline
% P16 (c1) & 186,82 & 4339,29 & 2,90 &  1,453 & -0,274 & -0,957 & \\
% P16 (c2) & 128,34 & 4339,29 & 2,90 &  1,453 & -0,274 & -0,970 &\\
% P16 (c3) & 190,03 & 4339,29 & 2,90 &  1,453 & -0,274 & -0,956 &\\
% \hline
% \end{tabular}

% \end{table}


% \begin{table}[H]
% \centering
% \caption{Verificação das tensões no solo}
% \label{tab:tensoes_solo}
% \setlength{\tabcolsep}{5pt}
% \begin{tabular}{lrrrrrrrrr}
% \hline
% El &
% $\sigma_{\max,c1}$ & $\sigma_{\min,c1}$ & $g_{c1}$ & $\sigma_{\max,c2}$ & $\sigma_{\min,c2}$ & $g_{c2}$ & $\sigma_{\max,c3}$ & $\sigma_{\min,c3}$ & $g_{c3}$\\
% \hline
% P16 & 43,99 & 41,21 & -0,934 & 35,56 & 22,97 & -0,946 & 49,63 & 37,03 & -0,925 \\
% \hline
% \end{tabular}
% \end{table}


% Após a determinação das tensões solicitantes para todas as combinações de carregamento e a avaliação dos respectivos valores de cada restrição $g$, é selecionada a combinação que representa a situação mais crítica. Essa combinação, por representar a pior situação em termos de solicitação dentre as analisadas, é utilizada para o processamento da função-objetivo.

% Neste caso, a partir da analise dos valores de $g_{rd,1}$ e $g_{rd,2}$ apresentados na tabela \ref{tab:puncao_secao_c_linha_p16} e \ref{tab:puncao_secao_c}, observa-se que a combinação $c2$ apresenta a condição mais crítica. Dessa forma, essa combinação é adotada para dar seguimento ao processo de otimização.

